% ============================================================
% DROP-IN REPLACEMENT for Draft_5c.tex lines 1179-1347
% (\subsection{Stability Properties} ... "held there by a boundary
%  condition.")
%
% Keeps labels : sec:stability, eq:kinematic_model, ass:trench,
%                eq:normal_form, eq:n_dynamics, thm:separatrix,
%                eq:iss_bound, prop:equivariance, eq:s1_lyapunov
% New labels   : eq:transverse_rates, tab:bandwidth
% Retired      : eq:oecs_n_dynamics, the z_k definition,
%                the tan(2*pi*n) footnote
% Note         : eq:normal_form now covers a generic surrogate Sigma
%                and uses kappa_perp in place of lambda_perp.
%                See diff D6 for the matching Appendix A edit.
% ============================================================

\subsection{Stability Properties}
\label{sec:stability}

The analysis is at the navigation level, with the cluster centroid a
kinematic integrator
\begin{equation}
    \dot{\mathbf{p}}_c = k\,\boldsymbol{\delta}(\mathbf{p}_c),
    \label{eq:kinematic_model}
\end{equation}
$\boldsymbol{\delta}$ the right-hand side of the active law,
(\ref{eq:d_law}) or (\ref{eq:oecs_law}), and $k > 0$ the navigation
gain. This abstracts the layers beneath, whose dynamics settle in
roughly 0.3~s against transit times of tens of seconds \cite{14}.
Estimates are taken exact, so hats are dropped; proofs are in
Appendix~\ref{app:proofs}.

\begin{assumption}[Trench and flow structure]
\label{ass:trench}
Let $\Sigma$ be the tracked surrogate, $D$ or $s_1$. On a tube $T_w =
\{|n| \leq w\}$ about a segment $\Gamma$, in arclength coordinates
$(\ell, n)$ along and across it,
\begin{equation}
    \Sigma(\ell, n) = h(\ell)
        + \tfrac{1}{2}\kappa_\perp(\ell)\,n^2,
    \qquad 0 < \underline{\kappa} \leq \kappa_\perp(\ell),
    \label{eq:normal_form}
\end{equation}
with vanishing cross terms. The curve $\Gamma$ is invariant under the
flow, with $\lVert\mathbf{v}\rVert \geq f_{\min} > 0$ outside balls of
radius $r_0$ around the flow saddles, and $|h'(\ell)| \geq m_h > 0$
wherever the band test fails on $\Gamma$.
\end{assumption}

The transverse gradient $\kappa_\perp n$ vanishes on $\Gamma$ and
nowhere else in the tube, so the problem splits: regulate $n$ to zero,
transport along $\ell$.

The two controllers are the same machine with one number changed. Both
command the transverse law
\begin{equation}
    \dot{n} = -c_{\max}\tanh\!\left(
        \frac{a_\perp\,n}{c_{\max}}\right),
    \label{eq:n_dynamics}
\end{equation}
and differ only in $a_\perp$. The $D$ tracker normalizes by the
estimated curvature, giving the Newton argument $-\kappa_\perp
n/\kappa_\perp = -n$; the $s_1$ tracker has no usable curvature
(Remark~\ref{rem:s1hess}) and substitutes the fixed gain $g_\perp$:
\begin{equation}
    a_{\perp,D} = 1,
    \qquad
    a_{\perp,s_1} = g_\perp\,\kappa_\perp(\ell).
    \label{eq:transverse_rates}
\end{equation}
One rate is fixed by design, the other set by the field.

Two regimes live inside (\ref{eq:n_dynamics}): far out $\tanh$ saturates
into a constant approach at $c_{\max}$, close in it is linear and the
command is proportional, with crossover at $|n| = c_{\max}/a_\perp$.
With $V = \tfrac{1}{2}n^2$, $\dot{V} = -c_{\max}\,n\tanh(a_\perp
n/c_{\max}) < 0$ for $n \neq 0$, so $T_w$ is forward invariant, the
proportional region is entered in finite time, and inside it concavity
of $\tanh$ gives $\dot{V} \leq -2a_\perp\tanh(1)V$.

Once $n \approx 0$ the tangential command takes over, with $\tau \geq 0$
in every branch of (\ref{eq:tau}), leaving one-dimensional transport in
the flow direction. The two phases do not interfere, and the reason is
visible in the law rather than in the analysis: \emph{the transverse
command contains no $\tau$}. The tangential mode switch is therefore
invisible to the regulated variable and cannot move the team across the
structure, and the $s_1$ tracker's projector enforces the same
separation explicitly.

\begin{theorem}[Network traversal]
\label{thm:separatrix}
Under Assumption~\ref{ass:trench}, with exact estimates, every
trajectory starting in $T_w$ (i) remains in $T_w$ and converges
exponentially to $\Gamma$, and (ii) traverses $\Gamma$ in the direction
of the ambient flow, reaching the $r_0$ neighborhood of a flow saddle in
finite time bounded by $L_\Gamma/m$, with $L_\Gamma$ the length of
$\Gamma$ and $m > 0$ the realized tangential speed.
\end{theorem}

At a flow saddle $D$ has a local minimum with $\mathbf{H}_D \succ 0$ and
the degenerate case of (\ref{eq:d_law}) is active. In practice the $D$
tracker continues past it: the floor $e$ of Section~\mbox{III-D} leaves
the estimated eigenvalues indefinite and an order of magnitude too small
there even without noise, so the tracker rides onto the crossing segment
and repeats (i) and (ii) with $\Gamma$ relabeled. Both behaviors are the
same theorem on different premises.

Termination carries its own certificate. When the ride flag clears the
projector becomes the identity and (\ref{eq:oecs_law}) is full gradient
descent on the surrogate itself,
\begin{equation}
    \dot{s}_1 = -c_{\max}\lVert\nabla s_1\rVert
    \tanh\!\left(\frac{g_\perp\lVert\nabla s_1\rVert}{c_{\max}}\right)
    \;\leq\; 0,
    \label{eq:s1_lyapunov}
\end{equation}
with equality only at a critical point: nothing holds the team in place,
it has arrived. This is the vanishing-gradient test
Section~\mbox{III-D} admits rather than the definiteness test it does
not.

Estimation error, and any residual cross term $\kappa_\perp'(\ell)\,n$,
enters the transverse channel as a bounded disturbance $|d| \leq
\delta$. Then $\dot{V} < 0$ outside an interval, and the structure is
practically stable with
\begin{equation}
    \limsup_{t\to\infty}|n(t)| \;\leq\;
    \frac{c_{\max}}{a_\perp}
    \operatorname{artanh}\!\left(\frac{\delta}{c_{\max}}\right)
    \;\approx\; \frac{\delta}{a_\perp},
    \qquad \delta \ll c_{\max},
    \label{eq:iss_bound}
\end{equation}
the disturbance divided by the local slope of the restoring law, which
is why $a_\perp$ matters. Here $\delta$ is the transverse gradient
error, placed at $\nu/\tilde{\rho}^{\,2}$ by Table~\ref{tab:budget}, so
tracking error degrades linearly in measurement noise: no control law
changes that exponent, and only $a_\perp$ scales the result.

\begin{proposition}[Frame equivariance of the closed loop]
\label{prop:equivariance}
Let an observer rotate as $\mathbf{p}_Q = \mathbf{Q}(t)\mathbf{p}$ and
let the $s_1$ tracker run with exact estimates on the correspondingly
transformed field. Then $\boldsymbol{\delta}_Q =
\mathbf{Q}\boldsymbol{\delta}$ and $\mathbf{p}_{c,Q}(t) =
\mathbf{Q}(t)\mathbf{p}_c(t)$ up to the frame-transport term: the
tracker follows the same material path in every frame. The single
exception is the degenerate guard, reachable only on the first tracking
step.
\end{proposition}

Every quantity (\ref{eq:oecs_law}) consumes is invariant or co-rotating,
the tangent selection (\ref{eq:tangent_select}) included. The $D$
tracker admits no counterpart: its surrogate carries the
$O(\mathrm{Ro})$ term of Section~\mbox{III-D} into everything its law
reads.

Two bandwidth limits sit underneath all of this. Explicit Euler at step
$\Delta t$ makes the near-structure map $n \mapsto n(1 - \Delta t\,k\,
a_\perp)$: monotone below $\Delta t\,k\,a_\perp = 1$, stable below $2$,
admitting a limit cycle above it. The kinematic abstraction separately
needs $k\,a_\perp\tau_{\mathrm{mom}} \ll 1$, since a proportional region
narrower than the plant can follow is not a proportional region. Both
read off $a_\perp$ alone, and at the operating point of
Table~\ref{tab:params} they bind for one tracker only.

\begin{table}[!t]
\centering
\caption{Discrete-time and actuator conditions at the double-gyre
operating point of Table~\ref{tab:params} ($\Delta t = 0.1$, $k = 3.0$,
$\tau_{\mathrm{mom}} = 0.28$, $g_\perp = 1.0$). For the $s_1$ tracker
$\kappa_\perp$ is the local transverse curvature of the structure,
quoted as recovered by the fit and as computed analytically at the same
point (Section~\mbox{VI-B-3}).}
\label{tab:bandwidth}
\footnotesize
\setlength{\tabcolsep}{4pt}
\begin{tabular}{@{}lcccl@{}}
\hline
 & $a_\perp$ & $\Delta t\,k\,a_\perp$ &
 $k\,a_\perp\tau_{\mathrm{mom}}$ & \\
\hline
$D$ tracker & $1.0$ & $0.30$ & $0.84$ & monotone \\
$s_1$, fitted $\kappa_\perp = 4.0$ & $4.0$ & $1.20$ & $3.4$ & oscillatory \\
$s_1$, analytic $\kappa_\perp = 6.9$ & $6.9$ & $2.07$ & $5.8$ & limit cycle \\
\hline
\end{tabular}
\end{table}

Two consequences follow. The $s_1$ tracker's larger zero-noise tracking
error is more likely a discretization effect than an estimation one,
since (\ref{eq:iss_bound}) predicts the opposite ranking: a larger
$a_\perp$ gives a \emph{tighter} ultimate bound under equal disturbance,
and only a limit cycle explains a larger error where that bound cannot.
And the loop is stable at the fitted curvature but not at the true one,
so the estimator's underestimate of $\kappa_\perp$ is carrying the
margin. Halving $\Delta t$ separates the two in one run, shrinking a
limit cycle and leaving an estimation error untouched, with the $D$
tracker indifferent either way at $\Delta t\,k\,a_\perp = 0.3$.
